\documentclass[10pt,a4paper]{amsart}
\usepackage{amsmath,amssymb,amsthm,graphicx}
% Retained minimal preamble for the Stacks Project chapter more-morphisms.tex (More on Morphisms).
% Verbatim excerpt of the chapter preamble: Xy-pic setup (native preamble lines 14-18) and the native
% theorem declarations with the mathematical macros (native preamble lines 154-217). Copied mechanically;
% nothing added, removed or redefined. The wrapper class and the amsmath/amssymb/amsthm/graphicx base
% packages are supplied by the assembler, so the class/externaldocument part of the native preamble is omitted.
% For commutative diagrams we use Xy-pic
\usepackage[all]{xy}

% We use 2cell for 2-commutative diagrams.
\xyoption{2cell}
% Theorem environments.
%
\theoremstyle{plain}
\newtheorem{theorem}[subsection]{Theorem}
\newtheorem{proposition}[subsection]{Proposition}
\newtheorem{lemma}[subsection]{Lemma}

\theoremstyle{definition}
\newtheorem{definition}[subsection]{Definition}
\newtheorem{example}[subsection]{Example}
\newtheorem{exercise}[subsection]{Exercise}
\newtheorem{situation}[subsection]{Situation}

\theoremstyle{remark}
\newtheorem{remark}[subsection]{Remark}
\newtheorem{remarks}[subsection]{Remarks}

\numberwithin{equation}{subsection}

% Macros
%
\def\lim{\mathop{\mathrm{lim}}\nolimits}
\def\colim{\mathop{\mathrm{colim}}\nolimits}
\def\Spec{\mathop{\mathrm{Spec}}}
\def\Hom{\mathop{\mathrm{Hom}}\nolimits}
\def\Ext{\mathop{\mathrm{Ext}}\nolimits}
\def\SheafHom{\mathop{\mathcal{H}\!\mathit{om}}\nolimits}
\def\SheafExt{\mathop{\mathcal{E}\!\mathit{xt}}\nolimits}
\def\Sch{\mathit{Sch}}
\def\Mor{\mathop{\mathrm{Mor}}\nolimits}
\def\Ob{\mathop{\mathrm{Ob}}\nolimits}
\def\Sh{\mathop{\mathit{Sh}}\nolimits}
\def\NL{\mathop{N\!L}\nolimits}
\def\CH{\mathop{\mathrm{CH}}\nolimits}
\def\proetale{{pro\text{-}\acute{e}tale}}
\def\etale{{\acute{e}tale}}
\def\QCoh{\mathit{QCoh}}
\def\Ker{\mathop{\mathrm{Ker}}}
\def\Im{\mathop{\mathrm{Im}}}
\def\Coker{\mathop{\mathrm{Coker}}}
\def\Coim{\mathop{\mathrm{Coim}}}

% Boxtimes
%
\DeclareMathSymbol{\boxtimes}{\mathbin}{AMSa}{"02}

%
% Macros for moduli stacks/spaces
%
\def\QCohstack{\mathcal{QC}\!\mathit{oh}}
\def\Cohstack{\mathcal{C}\!\mathit{oh}}
\def\Spacesstack{\mathcal{S}\!\mathit{paces}}
\def\Quotfunctor{\mathrm{Quot}}
\def\Hilbfunctor{\mathrm{Hilb}}
\def\Curvesstack{\mathcal{C}\!\mathit{urves}}
\def\Polarizedstack{\mathcal{P}\!\mathit{olarized}}
\def\Complexesstack{\mathcal{C}\!\mathit{omplexes}}
% \Pic is the operator that assigns to X its picard group, usage \Pic(X)
% \Picardstack_{X/B} denotes the Picard stack of X over B
% \Picardfunctor_{X/B} denotes the Picard functor of X over B
\def\Pic{\mathop{\mathrm{Pic}}\nolimits}
\def\Picardstack{\mathcal{P}\!\mathit{ic}}
\def\Picardfunctor{\mathrm{Pic}}
\def\Deformationcategory{\mathcal{D}\!\mathit{ef}}
\newenvironment{reference}{}{}
\newenvironment{slogan}{}{}
\newenvironment{history}{}{}
\usepackage{cleveref}
\crefname{originallemma}{original lemma}{original lemmas}
\Crefname{originallemma}{Original lemma}{Original lemmas}
\makeatletter
\newlabel{coherent-lemma-quasi-coherent-affine-cohomology-zero}{{01XB}{source}{}{sourcepointer-1}{}}
\newlabel{coherent-lemma-quasi-coherent-affine-cohomology-zero@cref}{{[originallemma][0][]01XB}{[1][0][]source}{}{}{}}
\makeatother
\crefname{originallemma}{original lemma}{original lemmas}
\Crefname{originallemma}{Original lemma}{Original lemmas}
\makeatletter
\newlabel{coherent-lemma-quasi-compact-h1-zero-covering}{{01XF}{source}{}{sourcepointer-2}{}}
\newlabel{coherent-lemma-quasi-compact-h1-zero-covering@cref}{{[originallemma][0][]01XF}{[1][0][]source}{}{}{}}
\makeatother
\crefname{originallemma}{original lemma}{original lemmas}
\Crefname{originallemma}{Original lemma}{Original lemmas}
\makeatletter
\newlabel{coherent-lemma-relative-affine-cohomology}{{089W}{source}{}{sourcepointer-3}{}}
\newlabel{coherent-lemma-relative-affine-cohomology@cref}{{[originallemma][0][]089W}{[1][0][]source}{}{}{}}
\makeatother
\crefname{originallemma}{original lemma}{original lemmas}
\Crefname{originallemma}{Original lemma}{Original lemmas}
\makeatletter
\newlabel{definition-thickening}{{04EX}{source}{}{sourcepointer-4}{}}
\newlabel{definition-thickening@cref}{{[originallemma][0][]04EX}{[1][0][]source}{}{}{}}
\makeatother
\crefname{originallemma}{original lemma}{original lemmas}
\Crefname{originallemma}{Original lemma}{Original lemmas}
\makeatletter
\newlabel{limits-proposition-affine}{{05YU}{source}{}{sourcepointer-5}{}}
\newlabel{limits-proposition-affine@cref}{{[originallemma][0][]05YU}{[1][0][]source}{}{}{}}
\makeatother
\crefname{originallemma}{original lemma}{original lemmas}
\Crefname{originallemma}{Original lemma}{Original lemmas}
\makeatletter
\newlabel{morphisms-lemma-i-star-equivalence}{{01QY}{source}{}{sourcepointer-6}{}}
\newlabel{morphisms-lemma-i-star-equivalence@cref}{{[originallemma][0][]01QY}{[1][0][]source}{}{}{}}
\makeatother
\title{Any thickening of an affine scheme is affine}
\author{The Stacks Project Authors}
\date{stacks-project:more-morphisms.tex:local-snapshot-20261011; GNU FDL 1.2 or later}
\begin{document}
\maketitle
\noindent\textit{Editorial note.} Statement span is lines 156-165 of sources/stacks-more-morphisms/pinned-original-source.tex; proof span is lines 167-206, the complete author proof (38 body lines).\par\medskip
\noindent\textit{Editorial note.} Navigation stubs: 6 cited-but-not-included labels are supplied as non-load-bearing pointers with their pinned tags where the snapshot tag index has them.\par\medskip
\noindent\textit{Editorial note.} extra\_packages is empty on purpose: the retained preamble already loads xy with options and a second declaration clashes in the assembler wrapper.\par\medskip
\noindent\textit{Editorial note.} Text is retained verbatim, including the author's own typos and placeholders.\par\medskip
\noindent\textit{Editorial note.} Pool card span [167,170] disagreed with the recomputed span and was ruled a line-number defect (Chief ruling 2026-10-13, LANE\_BRIEF §16.5/§16.6: real spans are computed from the \textbackslash{}label, the statement union the whole proof environments); card follows the recomputed span.\par\medskip
\noindent\textit{Editorial note.} The statement's reference block contains the chapter's only \textbackslash{}cite\{EGA1\} in this span; satisfied by the card-local bibliography.tex transcribed from the snapshot my.bib entry EGA1 (Institute des Hautes Etudes Scientifiques, Publications Mathematiques 4, 1960).\par\medskip
\noindent\textit{Editorial note.} Chapter opening: first lemma of the thickening series (definition-thickening stub carried with it).\par\medskip
\noindent\textit{Editorial note.} Pool-card span [167,170] covered only the pointer proof; the collected proof span 167-206 adds the substantive finite-order-thickening proof per the Lead ruling of 2026-10-13 (real spans recomputed from the \textbackslash{}label). The statement's reference block \textbackslash{}cite\{EGA1\} is satisfied by a card-local bibliography.tex transcribed from the snapshot my.bib EGA1 entry; bibliography\_backend none (thebibliography environment directly).\par\medskip
\noindent\textit{Allowed standard prerequisites (editorial).}\begin{itemize}
\item Explicit prerequisite boundary. Limits, proposition limits-proposition-affine; Cohomology of Schemes, lemmas coherent-lemma-quasi-compact-h1-zero-covering, coherent-lemma-relative-affine-cohomology, coherent-lemma-quasi-coherent-affine-cohomology-zero; Morphisms, lemma morphisms-lemma-i-star-equivalence; EGA1, Proposition 5.1.9 (cited in the statement's reference block).
\end{itemize}\medskip
\ifdefined\hypertarget\hypertarget{sourcepointer-1}{}\fi
\noindent\textit{Source reference.} Original lemma 01XB: sources/stacks-more-morphisms/pinned-original-source.tex, source lines 206--206. Cited by the retained statement or proof but not part of the collected spans; supplied only as a navigation pointer. The argument does not depend on its content.\par\medskip
\ifdefined\hypertarget\hypertarget{sourcepointer-2}{}\fi
\noindent\textit{Source reference.} Original lemma 01XF: sources/stacks-more-morphisms/pinned-original-source.tex, source lines 206--206. Cited by the retained statement or proof but not part of the collected spans; supplied only as a navigation pointer. The argument does not depend on its content.\par\medskip
\ifdefined\hypertarget\hypertarget{sourcepointer-3}{}\fi
\noindent\textit{Source reference.} Original lemma 089W: sources/stacks-more-morphisms/pinned-original-source.tex, source lines 206--206. Cited by the retained statement or proof but not part of the collected spans; supplied only as a navigation pointer. The argument does not depend on its content.\par\medskip
\ifdefined\hypertarget\hypertarget{sourcepointer-4}{}\fi
\noindent\textit{Source reference.} Original lemma 04EX: sources/stacks-more-morphisms/pinned-original-source.tex, source lines 36--63. Cited by the retained statement or proof but not part of the collected spans; supplied only as a navigation pointer. The argument does not depend on its content.\par\medskip
\ifdefined\hypertarget\hypertarget{sourcepointer-5}{}\fi
\noindent\textit{Source reference.} Original lemma 05YU: sources/stacks-more-morphisms/pinned-original-source.tex, source lines 206--206. Cited by the retained statement or proof but not part of the collected spans; supplied only as a navigation pointer. The argument does not depend on its content.\par\medskip
\ifdefined\hypertarget\hypertarget{sourcepointer-6}{}\fi
\noindent\textit{Source reference.} Original lemma 01QY: sources/stacks-more-morphisms/pinned-original-source.tex, source lines 206--206. Cited by the retained statement or proof but not part of the collected spans; supplied only as a navigation pointer. The argument does not depend on its content.\par\medskip
\begin{lemma}
\label{lemma-thickening-affine-scheme}
\begin{slogan}
Affineness is insensitive to thickenings
\end{slogan}
\begin{reference}
The case of a finite order thickening is \cite[Proposition 5.1.9]{EGA1}.
\end{reference}
Any thickening of an affine scheme is affine.
\end{lemma}

\begin{proof}
This is a special case of
Limits, Proposition \ref{limits-proposition-affine}.
\end{proof}

\begin{proof}[Proof for a finite order thickening]
Suppose that $X \subset X'$ is a finite order thickening with $X$ affine. Then
we may use Serre's criterion to prove $X'$ is affine. More precisely, we will
use Cohomology of Schemes, Lemma
\ref{coherent-lemma-quasi-compact-h1-zero-covering}. Let $\mathcal{F}$ be a
quasi-coherent $\mathcal{O}_{X'}$-module. It suffices to show that
$H^1(X', \mathcal{F}) = 0$. Denote $i : X \to X'$ the given closed immersion
and denote
$\mathcal{I} = \Ker(i^\sharp : \mathcal{O}_{X'} \to i_*\mathcal{O}_X)$.
By our discussion of finite order thickenings (following
Definition \ref{definition-thickening}) there exists an $n \geq 0$
and a filtration
$$
0 = \mathcal{F}_{n + 1} \subset \mathcal{F}_n \subset
\mathcal{F}_{n - 1} \subset \ldots \subset
\mathcal{F}_0 = \mathcal{F}
$$
by quasi-coherent submodules such that $\mathcal{F}_a/\mathcal{F}_{a + 1}$ is
annihilated by $\mathcal{I}$. Namely, we can take
$\mathcal{F}_a = \mathcal{I}^a\mathcal{F}$. Then
$\mathcal{F}_a/\mathcal{F}_{a + 1} = i_*\mathcal{G}_a$ for some quasi-coherent
$\mathcal{O}_X$-module $\mathcal{G}_a$, see Morphisms, Lemma
\ref{morphisms-lemma-i-star-equivalence}. We obtain
$$
H^1(X', \mathcal{F}_a/\mathcal{F}_{a + 1}) =
H^1(X', i_*\mathcal{G}_a) = H^1(X, \mathcal{G}_a) = 0
$$
The second equality comes from Cohomology of Schemes, Lemma
\ref{coherent-lemma-relative-affine-cohomology}
and the last equality from Cohomology of Schemes, Lemma
\ref{coherent-lemma-quasi-coherent-affine-cohomology-zero}.
Thus $\mathcal{F}$ has a finite filtration whose successive quotients
have vanishing first cohomology and it follows by a simple
induction argument that $H^1(X', \mathcal{F}) = 0$.
\end{proof}

\begin{thebibliography}{1}
\bibitem{EGA1}
Grothendieck, Alexander and Dieudonn{\'e}, Jean.
\newblock {\em {\'E}l{\'e}ments de g{\'e}om{\'e}trie alg{\'e}brique {I}}.
\newblock Institute des {H}autes {\'E}tudes {S}cientifiques.
\newblock Publications {M}ath{\'e}matiques, Vol.~4, 1960.
\end{thebibliography}
\end{document}
